\section{Qualidade dos registros e cuidados adotados}
\label{sec:qualidade}

Nenhuma base administrativa é perfeita, e parte do trabalho
exploratório é localizar exatamente onde estão os problemas para que
eles não contaminem as conclusões. A boa notícia: os problemas deste
painel são poucos, identificados e circunscritos.

\textbf{Ocupação acima de 100\%.} Na internação, 14,7\% das observações
(509) ultrapassam 100\%, e na UTI, 6,1\% (209). Em sua maioria isso é
informação verdadeira de sobrecarga (mais diárias que a capacidade
cadastrada comporta) e não deve ser apagado. O problema real está na
cauda extrema da UTI: \textbf{os 15 maiores valores pertencem todos a
2020 e 2021}, liderados pelo \textbf{Hospital Irmãos Penteado e Santa
Casa de Campinas (CNES 2022648)}, que registra 875\% de ocupação de UTI
em 2021, com 6.389 diárias para apenas 16 leitos cadastrados. A
explicação é operacional: na onda da pandemia os hospitais expandiram
leitos de UTI às pressas e o cadastro não acompanhou, produzindo
divisões por um denominador defasado. O cuidado adotado: analisar a
ocupação como razão positiva, tratar a UTI em duas partes (ter ou não
UTI ativa; e, entre os que têm, a intensidade de uso) e aparar a cauda
extrema (winsorização no percentil 99: 124\% na internação e 164\% na
UTI), com atenção redobrada ao biênio pandêmico.

\textbf{Denominadores frágeis na pandemia.} Ao excluir as internações
COVID, alguns hospitais poderiam ficar com pouquíssimas internações
restantes, tornando suas taxas instáveis. No painel isso ocorre em
\textbf{apenas 2 observações, ambas do Hospital Infante D.\ Henrique
(CNES 2097613, São José do Rio Preto)}: 13 saídas não COVID em 2020 e
somente 3 em 2021, ano em que o TMP calculado degenera para zero (o
único zero de TMP de todo o painel). As duas observações recebem
marcação de denominador frágil e ficam de fora das análises de
resultado.

\textbf{Ruído estatístico dos hospitais pequenos.} A
Figura~\ref{fig:funil} organiza essa discussão: cada ponto é um
hospital num ano, posicionado pelo volume de saídas (eixo horizontal) e
pela mortalidade (eixo vertical); as linhas mostram a faixa em que a
mortalidade deveria variar se as diferenças fossem puro acaso
estatístico, faixa que é larga para hospitais pequenos e estreita para
os grandes. O resultado: \textbf{cerca de 69\% das observações caem
fora até da faixa mais generosa} (1.119 acima e 1.269 abaixo). Isso diz
duas coisas ao mesmo tempo: as diferenças de mortalidade entre
hospitais são reais e estruturais, ligadas ao perfil assistencial, e
não flutuação de sorte; e, abaixo de mil saídas por ano, o acaso passa
a pesar de verdade, o que justifica a cautela extra com hospitais
pequenos.

\textbf{Valores extremos e completude.} Pela regra usual de detecção
(3 intervalos interquartílicos acima do terceiro quartil), há 11
observações extremas de mortalidade (acima de 18,9\%), 82 de custo
(acima de R\$~4.727), 206 de alta complexidade e 223 de TMP, estas
últimas concentradas nos hospitais de longa permanência já descritos,
que merecem tratamento em separado e não corte cego. Quanto à
completude, \textbf{o painel não tem nenhum valor faltante nos sete
indicadores} nem nas variáveis de estratificação; as únicas lacunas
estão em campos cadastrais (nome fantasia e especialização), sem
consequência para as análises.

\begin{figure}[H]
\centering
\includegraphics[width=.96\textwidth]{fig_ae_16_funnel_mortalidade.png}
\caption{Funil da mortalidade: cada ponto é um hospital num ano; a
linha horizontal é a mortalidade média da rede (5,4\%) e as linhas
curvas delimitam a faixa esperada por puro acaso estatístico (95\% e
99,8\%), conforme o volume de saídas.}
\label{fig:funil}
\end{figure}

\comoler{o funil é estreito à direita (hospitais grandes, onde o acaso
pesa pouco) e largo à esquerda (hospitais pequenos, onde uma morte a
mais ou a menos muda a taxa). Pontos fora das linhas curvas indicam
mortalidade que o acaso não explica: a nuvem inteira transbordando das
linhas mostra que as diferenças entre hospitais são estruturais, e a
região à esquerda de mil saídas mostra onde as taxas devem ser lidas
com reserva.}
